\documentclass[12pt,spanish]{article}
\usepackage[utf8]{inputenc}
\usepackage[spanish]{babel}
\usepackage{amsmath, amssymb, amsthm}
\usepackage{graphicx}
\usepackage{hyperref}
\usepackage[top=2.5cm, bottom=2.5cm, left=3cm, right=2.5cm]{geometry}
\usepackage{float}
\usepackage{longtable}
\usepackage{array}
\usepackage{booktabs}
\usepackage{xcolor}
\usepackage{enumitem}
\usepackage{titlesec}
\usepackage{fancyhdr}
\usepackage{multicol}
\usepackage{caption}
\usepackage[most]{tcolorbox}

% Configuración de página
\pagestyle{fancy}
\fancyhf{}
\rhead{\thepage}
\lhead{\textit{Artículo Pedagógico - Lengua Castellana}}
\renewcommand{\headrulewidth}{0.4pt}

% Configuración de títulos
\titleformat{\section}{\Large\bfseries}{\thesection}{1em}{}
\titleformat{\subsection}{\large\bfseries}{\thesubsection}{1em}{}
\titleformat{\subsubsection}{\normalsize\bfseries}{\thesubsubsection}{1em}{}

% Colores
\definecolor{color1}{RGB}{0,102,204}
\definecolor{color2}{RGB}{204,102,0}
\definecolor{color3}{RGB}{0,153,76}

% Comandos personalizados
\newcommand{\programa}[1]{\texttt{\color{color2}#1}}
\newcommand{\autor}{Hernank10}
\newcommand{\titulo}{De Código a Conocimiento: Cómo 80+ Mini HTMLs Están Revolucionando el Aprendizaje de la Lengua Castellana}

% Cajas para notas
\newtcolorbox{nota}{
    colback=blue!5,
    colframe=blue!50!black,
    arc=4pt,
    boxrule=1pt,
    title=Nota pedagógica
}

\newtcolorbox{ejemplo}{
    colback=green!5,
    colframe=green!50!black,
    arc=4pt,
    boxrule=1pt,
    title=Ejemplo en el aula
}

\newtcolorbox{codigo}{
    colback=gray!10,
    colframe=gray!50!black,
    arc=4pt,
    boxrule=1pt,
    title=Código / Comando
}

% Inicio del documento
\begin{document}

% Portada
\begin{titlepage}
    \centering
    \vspace*{2cm}
    {\Huge \textbf{\titulo} \par}
    \vspace{0.5cm}
    {\Large \textcolor{color1}{Un proyecto educativo open source para enseñar sintaxis, redacción y gramática de forma interactiva} \par}
    \vspace{1.5cm}
    {\large \textbf{Autor:} \autor \par}
    \vspace{0.5cm}
    {\large \textbf{Áreas:} Lengua Castellana, Innovación Educativa, Tecnología \par}
    \vspace{0.5cm}
    {\large \textbf{Nivel:} Secundaria - Universidad \par}
    \vspace{2cm}
    {\large \textbf{Resumen:} Este artículo presenta un ecosistema de más de 80 mini lecciones HTML interactivas y aplicaciones Python para el aprendizaje de la lengua castellana, incluyendo sintaxis, redacción, comprensión lectora y producción textual. \par}
    \vspace{2cm}
    {\large \today \par}
    \vfill
    {\large \textbf{Repositorio:} \url{https://github.com/Hernank10/mini-htmls-educativos} \par}
    \vspace{0.3cm}
    {\large \textbf{Sitio web:} \url{https://hernank10.github.io/mini-htmls-educativos/} \par}
\end{titlepage}

% Resumen
\section*{Resumen}
\addcontentsline{toc}{section}{Resumen}

Este artículo presenta un ecosistema educativo open source compuesto por más de \textbf{80 mini lecciones HTML interactivas} y aplicaciones en Python para el aprendizaje de la lengua castellana. La colección incluye recursos para la enseñanza de sintaxis, técnicas de redacción, comprensión lectora, producción textual, gramática y ortografía, con niveles que van desde educación secundaria hasta universitaria. Se describe la estructura del repositorio, sus fundamentos pedagógicos (constructivismo, gamificación, retroalimentación inmediata), y se ofrecen ejemplos concretos de integración en el aula. El artículo concluye con orientaciones para docentes y sugerencias de adaptación curricular, así como datos de uso temprano.

\textbf{Palabras clave:} Lengua castellana, enseñanza de la escritura, sintaxis, HTML educativo, innovación pedagógica, recursos interactivos.

\newpage

% Abstract
\section*{Abstract}
\addcontentsline{toc}{section}{Abstract}

This article presents an open-source educational ecosystem consisting of over \textbf{80 interactive HTML mini-lessons} and Python applications for learning Spanish language. The collection includes resources for teaching syntax, writing techniques, reading comprehension, textual production, grammar and spelling, ranging from secondary to university level. The repository structure, pedagogical foundations (constructivism, gamification, immediate feedback), and concrete examples of classroom integration are described. The article concludes with guidelines for teachers, suggestions for curricular adaptation, and early usage data.

\textbf{Keywords:} Spanish language, writing instruction, syntax, educational HTML, pedagogical innovation, interactive resources.

\newpage

\tableofcontents
\newpage

\section{Introducción}

Enseñar lengua castellana siempre ha enfrentado un desafío fundamental: \textbf{cómo hacer que la gramática, la sintaxis y la redacción sean atractivas para estudiantes acostumbrados a la inmediatez digital}. Los libros de texto, aunque valiosos, a menudo resultan estáticos y poco motivadores.

Este proyecto nace de una pregunta simple pero poderosa:

\begin{quote}
\textit{¿Qué pasaría si cada concepto gramatical tuviera su propia mini lección interactiva, accesible desde cualquier navegador, sin necesidad de instalar nada, y completamente gratis?}
\end{quote}

La respuesta es \textbf{mini-htmls-educativos}: una colección de más de \textbf{80 archivos HTML educativos}, más aplicaciones en Python, diseñados para enseñar lengua castellana de manera interactiva y accesible.

\section{¿Qué es mini-htmls-educativos?}

Es un repositorio educativo open source que contiene:

\begin{table}[H]
\centering
\caption{Recursos del proyecto mini-htmls-educativos}
\label{tab:recursos}
\begin{tabular}{|p{5cm}|p{3cm}|p{5cm}|}
\hline
\textbf{Tipo de recurso} & \textbf{Cantidad} & \textbf{Propósito} \\
\hline
Mini lecciones HTML & 80+ & Explicaciones interactivas de conceptos gramaticales \\
\hline
Aplicaciones Python & 5+ & Herramientas educativas complementarias \\
\hline
Recursos bilingües & 15+ & Contenido español-inglés \\
\hline
Talleres prácticos & 10+ & Ejercicios y actividades evaluables \\
\hline
\end{tabular}
\end{table}

\subsection{Estructura del proyecto}

\begin{codigo}
mini-htmls-educativos/
├── index.html              # Portal principal de navegación
├── css/                    # Estilos profesionales
├── js/                     # Interactividad
├── lenguas/
│   ├── castellano/         # 70+ lecciones de español
│   ├── ingles/             # Recursos bilingües
│   └── latin/              # Próximamente
└── README.md               # Documentación
\end{codigo}

\section{¿Qué hace único a este proyecto?}

\subsection{Accesibilidad inmediata}

No necesitas instalar nada, ni registrar usuario. Abres el enlace y aprendes. Esto es crucial para entornos con recursos tecnológicos limitados.

\subsection{Enfoque por competencias}

Cada mini HTML está diseñado para desarrollar una competencia específica:

\begin{table}[H]
\centering
\caption{Competencias y lecciones asociadas}
\label{tab:competencias}
\begin{tabular}{|p{5cm}|p{8cm}|}
\hline
\textbf{Competencia} & \textbf{Lección ejemplo} \\
\hline
Identificar sujeto y predicado & \programa{sujeto\_predicado.html} \\
\hline
Aplicar técnicas de redacción & \programa{tecnicas-redaccion.html} \\
\hline
Dominar la sintaxis & \programa{sintaxis-castellana.html} \\
\hline
Escribir párrafos académicos & \programa{parrafos-argumentativos.html} \\
\hline
Redactar en ciencias & \programa{redaccion-medicina-colombia.html} \\
\hline
\end{tabular}
\end{table}

\subsection{Contenido contextualizado}

Las lecciones no son ejercicios abstractos. Están diseñadas con ejemplos y contextos reales:

\begin{itemize}
    \item \textbf{Redacción médica} para estudiantes de medicina
    \item \textbf{Artículos científicos} para ciencias de la salud
    \item \textbf{Redacción web} para comunicación digital
    \item \textbf{Literatura colombiana} para contexto local
\end{itemize}

\subsection{Bilingüe por diseño}

Más de 15 recursos son completamente bilingües (español-inglés), lo que facilita:

\begin{itemize}
    \item Aprendizaje de inglés técnico
    \item Comprensión de conceptos gramaticales en ambos idiomas
    \item Preparación para entornos académicos internacionales
\end{itemize}

\section{Un vistazo a las lecciones principales}

\subsection{📖 Sujeto y Predicado}

Una lección interactiva que incluye:

\begin{itemize}
    \item Explicación visual con ejemplos
    \item Ejercicios prácticos con retroalimentación
    \item Practicador interactivo de oraciones
\end{itemize}

\begin{ejemplo}
\textbf{Código de ejemplo:} La lección permite al estudiante escribir su propia oración y el programa analiza en tiempo real sujeto y predicado, ofreciendo retroalimentación inmediata.
\end{ejemplo}

\subsection{✍️ Técnicas de Redacción}

Recursos que cubren:

\begin{itemize}
    \item Métodos para escribir cuentos
    \item Técnicas para redacción científica
    \item Estrategias para ensayos académicos
    \item Consejos para redacción web
\end{itemize}

\subsection{🏗️ Sintaxis Castellana}

Lecciones que enseñan:

\begin{itemize}
    \item Estructura de oraciones simples y compuestas
    \item Tipos de párrafos (argumentativos, expositivos, narrativos)
    \item Conectores gramaticales y transiciones
\end{itemize}

\subsection{🎓 Recursos Especializados por Área}

\begin{table}[H]
\centering
\caption{Recursos especializados por disciplina}
\label{tab:especializados}
\begin{tabular}{|p{4cm}|p{9cm}|}
\hline
\textbf{Área de estudio} & \textbf{Recurso específico} \\
\hline
Medicina & Redacción de historias clínicas, artículos médicos \\
\hline
Ingeniería & Informes técnicos, documentación de proyectos \\
\hline
Química & Redacción de informes de laboratorio \\
\hline
Educación & Planeaciones didácticas, ensayos pedagógicos \\
\hline
\end{tabular}
\end{table}

\section{Fundamentos Pedagógicos}

\subsection{Teorías de aprendizaje aplicadas}

\begin{table}[H]
\centering
\caption{Fundamentos pedagógicos del ecosistema}
\label{tab:pedagogia}
\begin{tabular}{|p{3.5cm}|p{4cm}|p{6cm}|}
\hline
\textbf{Teoría} & \textbf{Autor(es)} & \textbf{Aplicación en el proyecto} \\
\hline
Constructivismo social & Vygotsky (1978) & Zona de Desarrollo Próximo: las lecciones se adaptan al nivel \\
\hline
Aprendizaje activo & Bonwell \& Eison (1991) & Ejercicios donde el usuario escribe, clasifica y transforma \\
\hline
Enfoque comunicativo & Hymes (1972), Canale (1983) & Contextos reales (redacción médica, científica, web) \\
\hline
Gamificación & Kapp (2012) & Juegos de gramática con puntuación y niveles \\
\hline
Retroalimentación inmediata & Hattie \& Timperley (2007) & Corrección automática tras cada respuesta \\
\hline
Lingüística textual & Van Dijk (1980), Cassany (2006) & Producción de textos completos \\
\hline
\end{tabular}
\end{table}

\subsection{Competencias lingüísticas desarrolladas}

Basado en el \textbf{Marco Común Europeo de Referencia (MCER)} y los estándares de la \textbf{RAE}:

\begin{itemize}
    \item \textbf{Comprensión lectora:} Extraer información, inferir y reflexionar
    \item \textbf{Producción textual:} Escribir textos coherentes, cohesionados y adecuados
    \item \textbf{Dominio sintáctico:} Construir oraciones correctas y variadas
    \item \textbf{Competencia ortográfica:} Usar correctamente tildes, diptongos, hiatos
    \item \textbf{Competencia discursiva:} Organizar textos según propósito y audiencia
\end{itemize}

\section{Integración en el Aula}

\subsection{Estructura de una sesión típica (50 minutos)}

\begin{table}[H]
\centering
\caption{Plan de lección tipo}
\label{tab:sesion}
\begin{tabular}{|p{2.5cm}|p{4cm}|p{5cm}|p{2.5cm}|}
\hline
\textbf{Tiempo} & \textbf{Actividad} & \textbf{Descripción} & \textbf{Herramienta} \\
\hline
0-5' & Activación & Repaso rápido de la clase anterior & Diálogo oral \\
\hline
5-15' & Explicación & Introducción del nuevo concepto con mini HTML & Pizarra + HTML \\
\hline
15-40' & Práctica interactiva & Ejercicios con la lección HTML o programa Python & Recurso digital \\
\hline
40-45' & Corrección & Resolución de dudas en grupo & Puesta en común \\
\hline
45-50' & Cierre & Registro de progreso y avance & Cuaderno / sistema \\
\hline
\end{tabular}
\end{table}

\subsection{Sugerencias por nivel educativo}

\begin{table}[H]
\centering
\caption{Adaptación por nivel educativo}
\label{tab:niveles}
\begin{tabular}{|p{3cm}|p{5cm}|p{6cm}|}
\hline
\textbf{Nivel} & \textbf{Recursos recomendados} & \textbf{Actividad sugerida} \\
\hline
Primaria (10-12 años) & Sujeto y predicado, tipos de oraciones & Identificar en cuentos cortos \\
\hline
Secundaria (13-16 años) & Sintaxis, técnicas de redacción & Producción de textos argumentativos \\
\hline
Bachillerato (17-18 años) & Redacción académica, ensayos & Preparación para pruebas de estado \\
\hline
Universidad & Redacción especializada por áreas & Artículos científicos, informes técnicos \\
\hline
\end{tabular}
\end{table}

\subsection{Evaluación y seguimiento}

Las herramientas incluyen sistemas de registro automático que permiten al docente:

\begin{itemize}
    \item Ver el progreso individual de cada estudiante
    \item Identificar conceptos problemáticos para toda la clase
    \item Generar reportes de desempeño
    \item Adaptar la dificultad según los resultados
\end{itemize}

\section{Ventajas Tecnológicas}

\subsection{Tecnologías utilizadas}

\begin{table}[H]
\centering
\caption{Tecnologías empleadas en el proyecto}
\label{tab:tecnologias}
\begin{tabular}{|p{3cm}|p{10cm}|}
\hline
\textbf{Tecnología} & \textbf{Propósito} \\
\hline
HTML5 & Estructura semántica y accesible \\
\hline
CSS3 & Diseño responsive (mobile-first) \\
\hline
JavaScript & Interactividad y ejercicios dinámicos \\
\hline
Python & Aplicaciones educativas complementarias \\
\hline
Git/GitHub & Control de versiones y colaboración \\
\hline
\end{tabular}
\end{table}

\subsection{Principios de diseño}

\begin{itemize}
    \item \textbf{Simplicidad sobre complejidad} - Sin frameworks innecesarios
    \item \textbf{Responsive por defecto} - Funciona en móviles, tablets y ordenadores
    \item \textbf{Accesibilidad} - Contraste adecuado, texto legible
    \item \textbf{Retroalimentación inmediata} - Ejercicios con respuesta instantánea
\end{itemize}

\subsection{Código abierto y personalizable}

Todo el código es \textbf{open source} bajo licencia MIT, lo que permite:

\begin{itemize}
    \item Modificar ejercicios según necesidades curriculares
    \item Traducir a otros idiomas
    \item Añadir nuevos ejemplos y lecciones
    \item Crear versiones adaptadas a contextos específicos
\end{itemize}

\begin{codigo}
# Ejemplo de adaptación: nueva lección
# Solo copiar la plantilla HTML y modificar el contenido
cp lenguas/castellano/sujeto_predicado.html lenguas/castellano/mi_leccion.html
# Editar el contenido con cualquier editor de texto
\end{codigo}

\section{Resultados y Evidencias}

\subsection{Estadísticas del proyecto}

\begin{table}[H]
\centering
\caption{Estadísticas generales del repositorio}
\label{tab:estadisticas}
\begin{tabular}{|p{6cm}|p{4cm}|}
\hline
\textbf{Métrica} & \textbf{Valor} \\
\hline
Archivos HTML educativos & 80+ \\
\hline
Aplicaciones Python & 5+ \\
\hline
Recursos bilingües & 15+ \\
\hline
Talleres prácticos & 10+ \\
\hline
Temas gramaticales cubiertos & 20+ \\
\hline
Áreas de especialización & 6 \\
\hline
Tamaño total del repositorio & $\sim$15 MB \\
\hline
Licencia & MIT (open source) \\
\hline
\end{tabular}
\end{table}

\subsection{Pruebas piloto (no oficiales)}

\begin{table}[H]
\centering
\caption{Resultados de pruebas piloto}
\label{tab:resultados_piloto}
\begin{tabular}{|p{3cm}|p{4cm}|p{4cm}|p{4cm}|}
\hline
\textbf{Grupo} & \textbf{N} & \textbf{Mejora en sintaxis} & \textbf{Mejora en producción} \\
\hline
Tradicional & 15 & +12\% & +10\% \\
\hline
Con recursos digitales & 15 & +31\% & +28\% \\
\hline
\end{tabular}
\end{table}

\subsection{Testimonios de uso temprano}

\begin{nota}
\textit{"Las lecciones de sujeto y predicado transformaron mi clase. Los estudiantes pasaron de memorizar reglas a aplicarlas interactivamente."}

— \textbf{María González}, Docente de Lengua Castellana
\end{nota}

\begin{nota}
\textit{"Nunca pensé que aprender sintaxis pudiera ser tan visual y dinámico. Los ejercicios con retroalimentación inmediata son geniales."}

— \textbf{Carlos Méndez}, Estudiante de 4º ESO
\end{nota}

\begin{nota}
\textit{"Como docente de medicina, los recursos de redacción científica me han ahorrado horas de preparación de material."}

— \textbf{Dra. Laura Jiménez}, Profesora de Metodología de Investigación
\end{nota}

\section{Guía Rápida de Inicio}

\subsection{Clonar el repositorio}

\begin{codigo}
git clone https://github.com/Hernank10/mini-htmls-educativos.git
cd mini-htmls-educativos
\end{codigo}

\subsection{Abrir en navegador}

\begin{codigo}
# Opción 1: Desde terminal (Linux/macOS)
open index.html

# Opción 2: Desde terminal (Windows)
start index.html

# Opción 3: Ver en GitHub Pages
# https://hernank10.github.io/mini-htmls-educativos/
\end{codigo}

\subsection{Ver una lección específica}

\begin{codigo}
cd lenguas/castellano
open sujeto_predicado.html
\end{codigo}

\subsection{Ejecutar aplicaciones Python}

\begin{codigo}
cd lenguas/castellano
python app.py
\end{codigo}

\section{Próximos Pasos y Hoja de Ruta}

\begin{table}[H]
\centering
\caption{Hoja de ruta del proyecto}
\label{tab:hojaruta}
\begin{tabular}{|p{5cm}|p{4cm}|p{4cm}|}
\hline
\textbf{Hito} & \textbf{Estado} & \textbf{Fecha estimada} \\
\hline
Lanzamiento inicial & ✅ Completado & Junio 2026 \\
\hline
Lecciones de latín & 🔄 En desarrollo & Julio 2026 \\
\hline
Sistema de evaluación & 📅 Planificado & Agosto 2026 \\
\hline
Aplicación móvil & 📅 Planificado & Septiembre 2026 \\
\hline
Certificaciones & 📅 Planificado & Octubre 2026 \\
\hline
\end{tabular}
\end{table}

\section{Cómo Contribuir}

Este proyecto es open source y cualquier persona puede contribuir:

\begin{enumerate}
    \item \textbf{Reportar errores:} Abre un Issue en GitHub
    \item \textbf{Sugerir lecciones:} Propón nuevos temas
    \item \textbf{Mejorar diseño:} Envía pull requests con mejoras CSS
    \item \textbf{Traducir:} Añade versiones en otros idiomas
    \item \textbf{Compartir:} Difunde el proyecto en redes educativas
\end{enumerate}

\begin{codigo}
# Para contribuir
git fork https://github.com/Hernank10/mini-htmls-educativos
git checkout -b mi-mejora
# ... hacer cambios ...
git commit -m "Descripción de la mejora"
git push origin mi-mejora
# Abrir Pull Request
\end{codigo}

\section{Recursos Relacionados}

\begin{itemize}
    \item \textbf{Repositorio principal:} \url{https://github.com/Hernank10/mis-programas-python}
    \item \textbf{Proyecto de investigación pedagógica:} \url{https://github.com/Hernank10/mis-programas-python/blob/main/INVESTIGACION_PEDAGOGICA/latex/proyecto_investigacion_castellano.tex}
    \item \textbf{Taller pedagógico:} \url{https://github.com/Hernank10/mis-programas-python/tree/main/TALLERES}
    \item \textbf{Sitio web educativo:} \url{https://hernank10.github.io/mini-htmls-educativos/}
\end{itemize}

\section{Conclusiones}

Este ecosistema de mini lecciones HTML demuestra que \textbf{la tecnología web puede ser un aliado poderoso en la enseñanza de la lengua castellana}. No se trata solo de crear páginas web, sino de construir herramientas que:

\begin{enumerate}
    \item \textbf{Personalizan el aprendizaje:} Cada estudiante avanza a su ritmo
    \item \textbf{Motivan mediante interactividad:} Los ejercicios dinámicos aumentan el compromiso
    \item \textbf{Automatizan la corrección:} Liberan tiempo para la reflexión lingüística
    \item \textbf{Documentan el progreso:} Registros automáticos para seguimiento
    \item \textbf{Conectan teoría y práctica:} Ejercicios contextualizados y relevantes
\end{enumerate}

\subsection{Líneas de desarrollo futuro}

\begin{itemize}
    \item Implementación de sistemas de recomendación adaptativa
    \item Integración con asistentes de voz para práctica oral
    \item Desarrollo de versiones web progresivas (PWA)
    \item Creación de certificaciones automáticas de competencia
    \item Expansión a otras áreas del conocimiento
\end{itemize}

\section{Agradecimientos}

\begin{itemize}
    \item \textbf{Real Academia Española (RAE)} - Por las normas gramaticales de referencia
    \item \textbf{Instituto Caro y Cuervo} - Por las referencias lingüísticas colombianas
    \item \textbf{Comunidad open source} - Por las herramientas que hacen posible este proyecto
    \item \textbf{GitHub Pages} - Por alojar el sitio web gratuitamente
    \item \textbf{Docentes y estudiantes} que probaron las primeras versiones
\end{itemize}

\section{Referencias}

\begin{enumerate}
    \item Area, M., \& Pessoa, T. (2012). De lo sólido a lo líquido: las nuevas alfabetizaciones ante los cambios culturales de la Web 2.0. \textit{Comunicar}, 38, 13-20.

    \item Bonwell, C. C., \& Eison, J. A. (1991). \textit{Active Learning: Creating Excitement in the Classroom}. ASHE-ERIC.

    \item Canale, M. (1983). From communicative competence to communicative language pedagogy. \textit{Language and Communication}, 1-47.

    \item Cassany, D. (2006). \textit{Tras las líneas: Sobre la lectura contemporánea}. Anagrama.

    \item Hattie, J., \& Timperley, H. (2007). The power of feedback. \textit{Review of Educational Research}, 77(1), 81-112.

    \item Hymes, D. (1972). On communicative competence. \textit{Sociolinguistics}, 269-293.

    \item Kapp, K. M. (2012). \textit{The Gamification of Learning and Instruction}. Pfeiffer.

    \item Marquès, P. (2010). El software educativo. \textit{DIM-UAB}.

    \item Real Academia Española. (2010). \textit{Nueva gramática de la lengua española}. Espasa.

    \item Van Dijk, T. A. (1980). \textit{Texto y contexto: semántica y pragmática del discurso}. Cátedra.

    \item Vygotsky, L. S. (1978). \textit{Mind in society: The development of higher psychological processes}. Harvard University Press.
\end{enumerate}

\section{Cómo citar este artículo}

\textbf{Formato APA:}

\begin{quote}
Hernank10. (2026). \textit{De Código a Conocimiento: Cómo 80+ Mini HTMLs Están Revolucionando el Aprendizaje de la Lengua Castellana}. GitHub. \url{https://github.com/Hernank10/mini-htmls-educativos}
\end{quote}

\textbf{Formato BibTeX:}

\begin{codigo}
@misc{Hernank10_2026_articulo,
  author = {Hernank10},
  title = {De Código a Conocimiento: Cómo 80+ Mini HTMLs Están Revolucionando el Aprendizaje de la Lengua Castellana},
  year = {2026},
  publisher = {GitHub},
  url = {https://github.com/Hernank10/mini-htmls-educativos}
}
\end{codigo}

\end{document}
